\subsection{Direct irreducible-curve propagation}
\label{subsec:selmer-p1-direct}

Only the absolutely irreducible branch of ordinary propagation is
needed for the rank-four characteristic-prime argument. Here we
verify that branch separately. In particular it needs neither the
odd-family second-order estimates nor incidence at nonsemisimple
points.

Retain the ordinary determinant ring, its actual Hecke ideal $J$,
the CM extension $L/F$, the finite allowed ramification, and the
hypotheses of Theorem~\ref{thm:selmer-arithmetic-propagation}.
Write $\Gamma_0$ for the corresponding Galois group over $F$,
$G_0$ for its index-two subgroup over $L$, and $c$ for the chosen
complex conjugation. Let $f:B\to A=k[[t]]$ be a NiceI point, with
$k$ finite of odd characteristic $\ell$. Let $K=k((t))$ and
$N=\ker(\epsilon|_{G_0})$, with $\epsilon$ the full cyclotomic
character. Write $\chi$ for its reduction in $k^\times$.
In the ordinary setup $\chi$ is nontrivial on $G_0$ (in particular
the distinguished split prime supplies such an element).
Use the normalized integral representation $\rho$ and symmetric
pairing of the arithmetic-propagation theorem. The normalization
and any further splitting of eigenvalues are made by finite
injective coefficient extension and change of basis, preserving
$f$'s kernel.

\begin{lemma}[Direct integral trace detection]
\label{lem:selmer-p1-direct-detection}
There exist finitely many $\sigma_j\in N$ having separable spectra
and a constant $C$ such that, for every $m\geq1$, the kernel of
the detecting map on $H^1(\Gamma_0,W/t^mW)$ has at most $|k|^C$
elements. Here $W=M_n(A)$ has the polarized dual action
\[
 g:Z\mapsto\chi(g)\rho(g)Z\rho(g)^{-1}\quad(g\in G_0),
 \qquad c:Z\mapsto-\chi(c)Z^T=Z^T ,
\]
and the detecting coordinates are
$\operatorname{tr}(z(\sigma_j)\rho(\sigma_j))$.
The same assertion holds after a finite injective coefficient
extension, with its new constant.
\end{lemma}

\begin{proof}
First work over $K$. Suppose the traces vanish on all of $N$.
Since $\chi|_N=1$, the map
$\dot\rho(g)=z(g)\rho(g)$ on $N$ satisfies the matrix derivation
rule. Extend $\rho,\dot\rho$ linearly to $K[N]$. If $\rho(a)=0$,
then for every $b\in K[N]$,
\[
 0=\operatorname{tr}\dot\rho(ab)
   =\operatorname{tr}(\dot\rho(a)\rho(b)).
\]
Absolute irreducibility on $N$ gives $\rho(K[N])=M_n(K)$.
The matrix trace pairing is perfect in every characteristic:
test against matrix units. Hence $\dot\rho(a)=0$. It follows
that $\dot\rho$ induces a derivation of $M_n(K)$, which is inner.
For example, for a derivation $D$ put
$X=\sum_i D(E_{i1})E_{1i}$ and apply the Leibniz rule to matrix
units to obtain $D=[X,-]$. Thus $z|_N$ is a coboundary.

After subtracting that coboundary, the cocycle on $G_0$ is zero
on $N$. The cocycle rule for $ng=g(g^{-1}ng)$ shows that every
value commutes with $\rho(N)$, hence is scalar. It therefore
descends to the abelian quotient $G_0/N$, with action $\chi$.
Choose $s$ with $\chi(s)\ne1$. Commutativity gives
\[
 (\chi(s)-1)z(g)=(\chi(g)-1)z(s).
\]
So the remaining cocycle is a coboundary as well.
Restriction from $\Gamma_0$ to $G_0$ is injective on $H^1$ by
inflation--restriction, since the quotient has order two.
This proves detection of the whole polarized dual cohomology;
no assumption on $n\bmod\ell$ is used. Each coordinate is
well-defined on cohomology by cyclicity of trace.

Here are the integral details, including cocycles which do not
lift from a finite truncation. For $i=1,2$,
$H^i(\Gamma_0,W)$ is the inverse limit of the finite
$H^i(\Gamma_0,W/t^mW)$. Indeed the cochain transition maps are
surjective (lift finite-valued continuous cochains), and the
inverse systems of finite cohomology groups satisfy the
Mittag--Leffler condition. The finite arithmetic cohomology
theorems apply because the allowed ramification set includes
the primes above $\ell$. In particular these inverse limits
are compact and Hausdorff. The sequence for multiplication by
$t$ gives an injection
\[
 H^i(\Gamma_0,W)/tH^i(\Gamma_0,W)
       \hookrightarrow H^i(\Gamma_0,W/tW).
\]
Lift a finite basis of the image; successive subtraction and
division by $t$ express every element as a convergent
$A$-linear combination of those lifts. Convergence in the
inverse-limit topology follows since $t^mH^i$ maps to zero at
precision $m$. Thus both $H^1$ and $H^2$ are finite $A$-modules.
Every continuous Laurent cocycle is bounded on the compact
group, so inverting $t$ gives precisely the corresponding
continuous $K$-cohomology.

By density, vanishing of the detecting traces on separable-spectrum
elements of $N$ implies their vanishing on all of $N$. Finite
dimensionality therefore selects finitely many such elements
whose trace functionals inject on $H^1(\Gamma_0,W)\otimes_A K$.
The exact sequence
\[
 0\longrightarrow H^1(\Gamma_0,W)/t^m
 \longrightarrow H^1(\Gamma_0,W/t^mW)
 \longrightarrow H^2(\Gamma_0,W)[t^m]\longrightarrow0
\]
bounds the number of nonliftable cosets by the fixed cardinality
of the $t$-power torsion of $H^2$. On the liftable part, Smith
normal form of the detecting map on the free quotient of $H^1$,
and the finite torsion of $H^1$, bound the kernel by a fixed
power of $|k|$. This proves the uniform assertion. A finite
extension of the DVR is finite free: expanding in a basis
commutes with continuous cochains and the trace coordinates.
Hence rational detection persists, and the same integral argument
applies with the new coefficient uniformizer.
\end{proof}

\begin{theorem}[The irreducible part of ordinary propagation]
\label{thm:selmer-p1-direct}
For the original ordinary Hecke ideal $J$ and every NiceI point
$f$ as above, $J\subseteq\ker f$ implies $J\subseteq p$
for every prime $p\subseteq\ker f$.
\end{theorem}

\begin{proof}
We use the actual auxiliary-prime construction and the ordinary
Taylor--Wiles modules of
Theorem~\ref{thm:selmer-taylor-wiles}, not a patching assertion
assuming nonempty support. Apply Lemma~\ref{lem:selmer-p1-direct-detection}
and split the selected characteristic polynomials by a finite
coefficient extension. The arithmetic construction of
Theorem~\ref{thm:selmer-arithmetic-propagation} gives $q$ distinct
split auxiliary primes at each level $N_0$, with residue norms
$1\bmod\ell^{N_0}$ and Frobenius approaching the selected
$\sigma_j$. It retains the entire original local inertia
conditions. We check the only Selmer estimate needed in this
irreducible case.

Let $V_m=M_n(A/t^m)$ be the primal polarized adjoint
(with $c$ acting by negative transpose), paired by the matrix
trace and the finite-field trace of the coefficient of $t^{m-1}$
with $W_m=W/t^mW$. Put $\nu=n(n-1)/2$ and $d=[F:\mathbf Q]$.
At ordinary places use the original upper-triangular condition
modulo local boundaries, with its diagonal restriction on the
torsion subgroup. The ordinary local count gives a local
Selmer-to-invariants ratio at most
$|k|^{m(n+\nu)d_v+C_v}$:
each diagonal character has $d_v+1$ parameters before
subtracting the $n$ diagonal boundary coordinates in the count,
and the $\nu$ upper roots have the usual local Euler contribution,
with bounded defects from their noncritical ratios. More precisely
this is the one-block specialization of
Proposition~\ref{prop:selmer-first-global-bound}'s local
calculation, using the scalar count
\eqref{eq:selmer-ordinary-scalar-count}.
Old unramified factors are one. At the distinguished place
the dual invariant module has rank zero by N4; at an auxiliary
place it has rank $n$ by the distinct roots. Their finite
invariant orders are bounded by $|k|^{C_0}$ and
$|k|^{mn+C_j}$, respectively, by the finite linear equation
count (the auxiliary Frobenius agrees at any fixed precision
eventually). The global invariant rank is zero: a commuting
matrix is scalar and a scalar skew matrix is zero. At each
real place the fixed rank is exactly $\nu$, since the
polarization is symmetric and $2$ is invertible.

For the relaxed dual group keep only the auxiliary-local
boundary conditions. An associated cocycle vanishes on each new
inertia group, because the coefficient action there factors
through the base quotient. It therefore descends to the base
Galois group: the new inertia groups normally generate the
kernel of the restriction to the original ramification set.
Equivalently use the pro-$\ell$ quotient and its kernel as in
Proposition~\ref{prop:selmer-tower-detection}. At fixed $m$
there are finitely many base cocycles; for all sufficiently
large levels their values at the chosen Frobenius and those
at $\sigma_j$ agree simultaneously. The local boundary
condition makes their detecting traces zero. The preceding
lemma thus bounds the order of the relaxed dual group by a
constant independent of $m$. Finite Poitou--Tate duality now
gives the first Selmer bound
\[
 \#\operatorname{Sel}_{N_0,m}
        \le |k|^{mn(d+q)+C},
\]
where the Selmer group is taken modulo polarized boundaries.
The exponent is
$(n+\nu)\sum_{v\mid\ell}d_v-\nu d+nq=n(d+q)$.
The constants may depend on this fixed point and the finite
auxiliary choice, but not on $m$ or sufficiently large $N_0$.

Pass to the one limiting group of the marked construction.
At each fixed precision the cocycle sets, local restrictions
and boundaries agree eventually: the finitely generated
profinite cover, the common limiting kernels and maps, and
the finite coefficient module give exactly the comparison
of Proposition~\ref{prop:selmer-limit-comparison}.
If $a$ bounded Laurent Selmer classes are independent over
$K$, choose integral representatives. Their independent
classes modulo boundaries give at least $|k|^{a(m-b)}$
classes at precision $m$, for one fixed $b$: choose finitely
many generator values testing the relation with a skew
primitive and put the resulting finite matrix in Smith
normal form. Thus $a\le n(d+q)=D_0$.
No second-order or odd-family statement is required.

For clarity we spell out the final dimension argument. Make
the harmless unramified coefficient extension and adjoin a
marked variable $s$ sent to the curve uniformizer, exactly
as in Theorem~\ref{thm:selmer-marked-propagation}. Form its
root-splitting algebra, limiting universal ring $R_0$, exact
relation ideal $I$, and map to the splitting algebra with
the diamond variables $Y$ killed. Use the same
discriminant multiplier, ordinary module and coefficient
variables. The construction of (MP2)--(MP9) uses no
second-order bound. Proposition~\ref{prop:selmer-patched-support}
therefore applies at the selected curve prime. Denote the
completed relation ring there by $R$ and its patched support
ideal by $S_R$. The latter is nonempty: the given
$J\subseteq\ker f$ and the exact $Y=0$ fibre identity give
support at the closed curve point. Every minimal support
prime $Z$ satisfies
\[
                     \dim R/Z\ge1+D_0 .
\]
The coefficient DVR is
$O'=\widehat{O[[s]]_{(\pi)}}$, with residue field $K$.

Since the residual determinant at this completed local
ring is absolutely irreducible, the one-block polarized
model of Proposition~\ref{prop:selmer-polarized-model}
is the full matrix algebra over $R$. Equivalently,
for the tangent test use Chenevier's absolutely irreducible
determinant representation theorem over $K[\delta]/(\delta^2)$,
and normalize the lifted symmetric pairing using $2^{-1}$.
The lifted pairing is symmetric: transposition changes the
intertwiner by a scalar of square one, reducing to $+1$,
hence equal to one. A symmetric pairing $1+\delta H$ is
made the original one by $1-\delta H/2$.
An $O'$-constant derivation of $R$ therefore gives a
polarized primal cocycle. Its actual ordinary identities
give the local flags and torsion conditions; the old
inertial and auxiliary equations give the remaining
Selmer conditions. All values have bounded Laurent
coordinates: the ring modulo the square of the curve
kernel is finite over $O[[s]]$, since its first
conormal module is finite over $k[[s]]$. This is the
first-order boundedness argument of
Proposition~\ref{prop:selmer-local-propagation}.

If its cocycle is a polarized boundary, the determinant
does not vary. Distinct ordinary characters and auxiliary
roots determine their lifted values by the ordered
ordinary equations and the separable polynomials.
The partial-root relations then force all diamond
coordinates to be constant, since the base auxiliary
inertia acts trivially. The limiting ring is generated
by these data and the marked coefficient ring; thus the
derivation is zero. Consequently
\[
 \dim_K\operatorname{Der}_{O'}(R,K)\le D_0 .
\]
Choose lifts of a basis of the relative cotangent space.
Completeness gives a surjection
$O'[[X_1,\ldots,X_{D_0}]]\to R$ (add zero generators
if necessary). The inverse image of any minimal support
prime $Z$ is a prime whose quotient has dimension at
least $1+D_0$. In this regular local domain of dimension
$1+D_0$ that prime must be zero. In particular the
kernel is zero, $R$ is a domain, and $S_R=0$.

Now every prime below the curve prime belongs to the
patched support after passing to the local completion.
The exact diamond-zero fibre identity gives containment
of $J$ at every such prime of the splitting algebra
over the original ring. Use flat going down and integral
lying over for the finite free splitting algebra and
unramified coefficient extension, and faithful flatness
of completion. Contracting the lifted primes (and
forgetting $s$) proves $J\subseteq p$ for every
$p\subseteq\ker f$, as asserted.
\end{proof}

This proves the actual NiceI conclusion directly from the
ordinary modules and first-order arithmetic estimate. The
general NiceE chain theorem remains available for the
other applications; the RH rank-four use of P1core need
not use its second-order presentation.
